\begin{itemize}
    \item As medições de bioimpedância são sensíveis a fatores fisiológicos e ambientais,
    como hidratação, alimentação recente e atividade física. O \sintonia\ trata esses
    fatores diretamente: registra-os na anamnese e os desconta com o modelo de correção
    \emph{within-user} (Seção~\ref{sec:correcao}), o que é um dos diferenciais do trabalho.
    Permanecem como limitação os fatores que a anamnese não captura, como o posicionamento
    do usuário e a temperatura da pele, e o fato de a correção reduzir o viés sem
    eliminá-lo, ainda sem validação contra DXA.

    \item O funcionamento da plataforma depende da comunicação Bluetooth Low Energy (BLE)
    entre o navegador e a balança, da disponibilidade da API Web Bluetooth e do protocolo
    BLE do dispositivo; mudanças de \emph{firmware} do fabricante que alterem esse protocolo
    podem exigir adaptações. Por outro lado, por ler a balança diretamente por BLE, o
    \sintonia\ não depende dos serviços em nuvem nem da infraestrutura do fabricante
    (Lefu/Unique Health), ao contrário do aplicativo original.

    \item A implementação atual foi desenvolvida e validada com a balança CF577BLE. A
    compatibilidade com outros modelos ou fabricantes não foi avaliada nesta etapa.

    \item A comunicação direta entre a aplicação web e dispositivos Bluetooth depende da API
    Web Bluetooth, atualmente suportada apenas por navegadores compatíveis, como Google
    Chrome e Microsoft Edge, o que pode limitar o uso em alguns ambientes.

    \item Os testes desta etapa usaram um número reduzido de medições, feitas pelos próprios
    autores, o que limita a generalização dos resultados para diferentes perfis
    populacionais.

    \item No estudo de campo previsto (Capítulo~\ref{cap:trabalhosfuturos}), pode haver viés
    de autosseleção, já que a participação dependerá do interesse espontâneo dos usuários.

    \item O trabalho não inclui medições próprias por absorciometria por dupla emissão de
    raios X (DXA), referência para avaliação da composição corporal. Por isso, as
    comparações se baseiam em modelos e equações da literatura e não constituem validação
    clínica direta da solução.

    \item A estimativa da impedância equivalente em 50~kHz é feita por interpolação
    log-linear a partir das medições em 20~kHz e 100~kHz. Embora compatibilize os dados da
    CF577BLE com modelos desenvolvidos para 50~kHz, é uma aproximação matemática sujeita a
    incertezas adicionais.

    \item Os mecanismos de correção e contextualização das medições foram pensados para
    reduzir variações ligadas a hidratação, alimentação e atividade física. Nesta etapa,
    porém, têm caráter exploratório e ainda precisam de validação com conjuntos de dados
    maiores e acompanhamento longitudinal.

    \item A correção \emph{within-user} usa as medições em condições ideais do próprio
    usuário como linha de base, o que assume que a composição corporal se mantém relativamente
    estável no período. Como o corpo muda ao longo do tempo (ganho ou perda de peso e de massa
    muscular ao longo de semanas ou meses), uma linha de base formada sobre um intervalo longo
    pode misturar estados fisiológicos diferentes e confundir a variação real da composição com
    o viés dos fatores de contexto. O número de medições e a janela de tempo usados para formar
    essa base ainda não são ajustados de forma adaptativa a essas mudanças.

    \item Embora o sistema use um SLM para gerar as interpretações, não foram realizadas nesta
    etapa técnicas de \emph{fine-tuning} supervisionado ou treinamento do modelo com bases
    próprias. A adaptação ao domínio foi feita por engenharia de \emph{prompt} (um \emph{system
    prompt} com instruções e \emph{guardrails}) e por regras de interpretação fundamentadas na
    literatura.

    \item Da mesma forma, não foi feita a utilização efetiva de bases sintéticas para
    treinamento ou ajuste do modelo. Estratégias com dados sintéticos poderão ser
    investigadas em estudos futuros.

    \item As interpretações da assistente se baseiam na literatura científica, nas regras de
    interpretação e nos modelos biométricos definidos, têm caráter informativo e educacional
    e não substituem avaliação clínica de profissionais habilitados.
\end{itemize}
